\chapter{黑格尔对自然法的批判}

在1802-1803年发表于《批判的哲学杂志》的论文中，黑格尔阐述了他对迄今为止的“自然法的科学探讨方式”的第一次大批判。在黑格尔学界，这篇论文在个别段落上争议极大，但在整体上却获得了近乎一致的评价。这种评价大体上概括如下：如果我们忽略这篇论文里面的一些晦涩之处，则可以认为，这篇自然法论文已经包含了1821年《法哲学原理》的纲要。如同在其他地方一样，这里，这种评价的文献来源是罗森克兰茨的《黑格尔传》一书。在罗森克兰茨看来，黑格尔在其成熟时期的法哲学中只不过是在一个精妙的体系中更为专门、详细地展示了所有的概念，相比之下，其构思的原创性已经在将其表达得“更美、更新鲜、部分地也更为真实”的杂志论文的“青年形态”中呈现出来。\footnote{罗森克兰茨，《黑格尔传》，柏林，1844年，第173页以下。}罗森克兰茨的判断在代表了1910年到1945年之间的新黑格尔主义的海林(Theodor Haering)和格洛克纳(Hermann Glockner)的黑格尔学述中延续了下来。格洛克纳在这一点上走得最远。他精炼地指出，黑格尔在其辩证方法最终形成之前，即已凭借自然法论文，作为伦理学家“完全”出现在我们面前了。\footnote{格洛克纳(H. Glockner)：《黑格尔》(Hegel)第2卷，第2版，斯图加特，1958年，第302页以下。也见海林(Th. Haering)：《黑格尔：他的意愿和他的著作》(Hegel. Sein Wollen und sein Werk)第2卷，莱比锡与柏林，1938年，第389页以下、404页以下。}

我们想问的是：情况真是这样吗？毫无疑问，就材料的丰富及其相对清晰的体系安排而言，自然法论文超出了青年黑格尔所有以前的著作。然而争议在于，黑格尔在基本概念的处理上是否有所超越，有多大程度的超越。由于对“绝对伦理”的那些概念（人民、政府等）的内容上的兴趣，新黑格尔主义跳过了这个问题。为了能够对这一问题做出决定，就需要反思为黑格尔自然法的伦理—政治范畴奠定基础之理论上的基本概念，首先就是“自然”概念。由于黑格尔自己认为需要对——正如他在《自然法》这一标题中表达的——自然和法的联系进行反思，因此之前的黑格尔学界居然没有提出这个本身就很明显的问题，就更是令人惊讶了。对这一联系进行考察，我们可以得知黑格尔法哲学的发展历程，以及黑格尔法哲学在近代自然法历史上所占据的地位。下面将会表明，这两个方面都依赖于黑格尔对其哲学本体论基础所获得的洞见。只有当辩证方法的论证已经进展到能够在自然法领域中为解决与传统自然法概念的争论提供指南时，黑格尔才作为伦理学家“完全”出现在我们面前。

\section*{1}

自然法论文在其批判的部分限于分析近代自然法，这一点对于论文的体系价值和历史价值来说具有核心的意义。在黑格尔看来，自然法在其历史发展过程中形成了两种彼此对立的科学探讨方式，也就是17世纪的“经验主义的探讨方式”和18世纪晚期的“形式主义的探讨方式”。前者从它认为是构成了实践经验之本质的众多自然关系和伦理关系中分离出一些个别因素（诸如自我保存的冲动、社交性、不合群性等），赋予它们概念统一性的形式，然后最终把它们提升为它所追求的科学体系的基本原理；后者则从经验中分离出统一性形式（对它来说，这种统一性形式乃是实践理性的本质），将它在自为的纯粹意志概念中固定下来，并把这种意志概念与杂多的经验关系对立起来。尽管两者间也存在着一种特别的区分，也就是说，其一的原则“是经验直观与普遍物的关系和混合，而另一种方式的原则则是绝对对立和绝对的普遍性”，但在黑格尔看来，这两种体系的“成分”（直观和概念）及其探讨方式（对规定性进行分离和固定，这些规定性的对立或不完善的连接），都是一样的。\footnote{参见《论自然法的科学探讨方式》，第328页以下、332页。黑格尔的这一著作以及《信仰与知识》和《费希特与谢林哲学体系的差别》等著作的页码，请参见《全集》第1卷，柏林，1832年。[这也包含了《全集》（百年纪念版）第1卷，H. 格洛克纳出版，斯图加特，1927年。因为这个版本也同时包含了第1版的页码。]}两者的共同基础是道德法则(lex moralis sive naturalis)与经验自然的分离，“自然法的概念及其在整个伦理科学中的处境的最新规定，均依赖于”这种分离。\footnote{同上书，第359页。}

黑格尔以17世纪的“经验主义的”探讨方式为例来阐明这一点。由于在纯粹经验中，自然的每一个因素似乎与其他因素一样，都是原初的，因此在追求使实践经验服从于概念统一性的过程中，纯粹经验就在方法论上被消解掉了，否定性的自然概念、“混沌”遂上升为体系的出发点。取代肯定性的、由“纯粹”经验把握到的自然（在这种自然中，“杂多之物的地位”并未遭到动摇）的，是一种否定性的自然概念。它要么在“存在的样态下被想象”为自然状态，要么在“可能性和抽象的形式”下表现为对于人类学上发现的能力、人的自然的列举。\footnote{《论自然法的科学探讨方式》，第333页以下。}这些规定、自然状态和人的自然的否定统一性是从伦理状态的肯定统一性中抽象出来的，然而后者并不能从前者推导出来。因此，经验主义的探讨方式就必须预设其他的自然规定性（比如社交的冲动和对暴死的恐惧），作为[从自然状态向社会状态]过渡的基础，或是寻求实定的、历史的规定性（比如强者对弱者的征服）的帮助。\footnote{同上书，第337页。}

对黑格尔来说，标示了17世纪自然法特征的体系基础，即对自然的否定，在康德—费希特哲学中臻于完成。一方面，黑格尔在那里承认，伦理总体在霍布斯、斯宾诺莎、普芬道夫的体系中达到了一种即便只是“没有形式的、外在的和谐”，另一方面，黑格尔在这里意在表明，观念论自然法的形式主义如何“徒劳地想要达到一个肯定的有机体”。\footnote{参见上书，第329、337页。}使得经验主义的探讨方式回溯到自然物的个别规定性成为可能的自然和概念的不完全分离，让位于一种完全的分离、“绝对的”对立；在[康德、费希特的]观念论体系中，自然显现为[作为]杂多的无本质的抽象物，而与一的同样的无本质的抽象物（也就是实践理性）相对立。\footnote{《论自然法的科学探讨方式》，第345页。}黑格尔在《费希特与谢林哲学体系的差别》（1801年）中就已经指出了从这种开端出发对于自然法原则和“人类共同体”建构必然得出的结论。自然法论文没有重复他在《差别》中得到的这种结论，而是以之为前提。就其实质而言，这种结论构成了批判自然法的经验主义探讨方式的反面。在黑格尔看来，在康德那里，特别是在费希特那里，自然就其依赖于概念而言只具有一种意义。费希特把这种依赖性提到概念支配自然的极端，但是通过将自我的这个自身无意识的产物本身\footnote{即自然。——译者}理解为自我，这种依赖性也并无改变。在费希特的自然法中，只是为了解释自我的“被规定的自然”，即冲动的自然[本性]，自然才被作为外在的自然（“规定性”）演绎出来。\footnote{参见《费希特与谢林哲学体系的差别》，第226页以下。黑格尔批判中的自然法概念的特殊价值来自于其哲学上的出发点，也就是从谢林那里接受过来的先验直观的客观设定(Objektivsetzen)方法。“由于坚持先验直观的主观性，自我仍然是一个主观的主体客体，这最鲜明地表现在自我与自然的关系上：一方面表现在自然的演绎上，另一方面表现在由此建立起来的许多科学上。”（第226页）（中译文参见黑格尔：《费希特与谢林哲学体系的差别》，宋祖良、程志民译，商务印书馆1994年版，第55页。——译者）}黑格尔如此总结他的批判：“在自然法中，理智的行动只是把自然作为一种可塑的质料产生出来，因此这就不是自由的、观念的行动，不是理性的行动，而是知性的行动。”\footnote{出处同上，第247页以下；也见《信仰与知识》，第140页以下。}

黑格尔的自然法概念试图重建法和自然的原始关系，从而与由费希特体系中“特别激烈地”表达出来的概念所造成的对自然的肢解相对立。\footnote{参见《费希特和谢林哲学体系的差别》，第234页：“在《自然法体系》所做的对自然的演绎和陈述中，自然和理性的绝对对立[……]表现得特别激烈。”如同自然是“由自我所设定的被设定存在”（第226页），自由因而也就是一个“纯然的否定物”（第237页）、绝对无规定性，其规定只有在对其自身不断地进行限制的情况下才是可能的，——“而限制的概念构成自由的王国，在这个自由的王国里，由于生命物被撕裂为概念和物质，由此，生命的每一个真正自由的、自为无限和不受限制的，也就是美好的相互关系就被毁灭了，自然就成了被统治的东西”（第236页以下）。黑格尔一再强调，费希特并不把自然理解为活的、自立的、发挥作用的东西，理解为“物体的内在充盈和力量的表达”（第248页）。参见第230页：“因此，不仅在理论方面，而且在实践方面，自然都是一个本质上被规定的东西，是一个死物”；第232页：理性在费希特那里不过是“自我的无依赖性以及自然的绝对被规定存在的理念，自然被设定为一个要被否定的东西、绝对依赖的”；第233页：“先验的观点[……]从显现为自然的东西中抽出了主体—客观性，只为自然保留了客观性的死壳”；第235页：“自然的基本特征”是“成为绝对对立的东西”。参见《信仰与知识》，第136页以下、140、142页。}黑格尔认为自然法的任务在于，“按照其名称”建构起“真正的肯定物”，——“自然法应该建构出名副其实的伦理自然”。\footnote{《论自然法的科学探讨方式》，第397页。格洛克纳没有注意到这个地方，他错误地主张，黑格尔之所以在自然法论文中把法哲学称为自然法科学，是因为在18世纪时并没有其他的名称，不是“因为他自己支持今天我们称之为自然法的思维方式的那种立场”。参见格洛克纳：《黑格尔》第2卷，第304页。}自然法名称中所包含的肯定物，导致黑格尔走向一种绝对设定自然和伦理的方法，这一方法是黑格尔法哲学的第一个发展阶段所特有的。\footnote{参见《论自然法的科学探讨方式》，第372页：为了能够规定自然法的真实概念及其与实践科学的关系，就必须突出“无限性方面”。也就是说，要预设“肯定的东西”即“绝对伦理总体就是一族人民”。进一步参见《伦理体系》，第415、406页以下、462页（页码根据黑格尔：《政治学与法哲学文集》，G. 拉松出版，莱比锡，1913年）；《差别》，第242页。}也就是说，黑格尔在与奠定近代自然法体系之基础的否定性的自然概念发生争论时，不得不重新回到近代自然法体系以前的法和社会的自然理论上。这就是黑格尔在术语和思想上与谢林发生关联的意义之所在，我们可以在他耶拿早期的作品中看到这种关联。黑格尔从谢林的自然哲学借来的自然概念，是传统的目的论自然概念，黑格尔卓有成效地将其重新用来为自然法奠基。\footnote{参见《信仰与知识》，第141、148页。黑格尔在这里把“服从自然和神圣必然性的永恒规律”称作真正伦理的原则。在这里，自然就是“活生生的东西自身，它在规律中同时设定自身是普遍的，并在人民中成为真正客观的”（第149页）。参见《伦理体系》，第460页以下。}这种自然理论的独特之处在于，自然不是以个别规定性的形式，而是作为整体自然进入伦理体系的；它既不与一个无论怎样理解的绝对概念（自我、自由、理性）相对立，也不作为否定性的自然状态理论与法权状态的实定权力相对立。它不与前者对立，是因为“伦理自然”在自身中统一了（纯粹个别性和普遍性的）概念；\footnote{《论自然法的科学探讨方式》，第395页。}不与后者对立，是因为伦理自然包含了“作为彻底同一的自然状态和威严”。而法权状态的威严对黑格尔而言，“自身不过就是绝对的伦理自然”。\footnote{同上书，第338页。}与古典自然法传统相一致，自然规律和道德法则彼此直接地关联在一起，公民状态被重新视为自然存在的实现，而非其克服。黑格尔在影射霍布斯及其后继者的自然状态理论时说道：“然而，那种在伦理关系中必须认为是要予以抛弃的自然物，本身就不是伦理，因而在其本来的形式中也就完全不能表现伦理。”\footnote{同上。除此之外，这个在原则上为目的论的自然概念还促使把前述“伦理自然”与（理性或自然的）“有机体”的概念等同起来，并将这个概念与“机械”的概念、自然法的“知性国家”进行对比。参见《费希特与谢林哲学体系的差别》，第242页。}

黑格尔对近代自然法的批判恰恰就是从这一点出发的：思维一种自然概念，这种自然概念不是一种“要予以抛弃的东西”，而是自身就是“伦理的”，是“伦理自然”。\footnote{这一措辞在《伦理体系》和自然法论文中一再出现，而且在某种程度上，可以认为它对黑格尔早期的自然法观念具有建构性意义。参见《伦理体系》，第443、471、473、478页；《论自然法的科学探讨方式》，第338、347页以下、369、393、397页。}按照黑格尔的观点，把伦理事物思维为“自然”，也就意味着：思维其一切潜能阶次的持存，思维这些潜能阶次的实在性以及它们与必然性的统一。因为“第一位的东西不是个体的个别性，而是伦理自然的生命(Lebendigkeit)、神性，对于伦理自然的本质来说，个别的个体太贫乏了，无法在其全部实在性中把握它的自然。”\footnote{《伦理体系》，第470页以下。对于黑格尔而言，个体只能“短暂地”（也就是说在时间中）、个别地表达伦理事物，作为否定物的个体通过时间的否定而被设定起来，以至于伦理事物的统一性（“无差异”）在个体中只不过是“形式的”而不是“实在的”；由此可以得出它（指伦理事物的统一性）与自然的等同：“[……]伦理事物必须把自身把握为自然，把握为一切潜能阶次的持存，并在其生动形态中把握每一潜能阶次，它必然与必然性合一，并作为相对同一性而持存。”（第471页）}这个具有多重含义的措辞\footnote{当指“自然”。——译者}应当把伦理整体对其部分、众多个体的优先性以及构成伦理整体一切环节的独立持存都表达出来。显然，这里自然概念用于两个方面：(1)作为总体，在斯宾诺莎“神或自然”(deus sive natura)的意义上，和(2)作为“本质”，在亚里士多德的城邦“按照自然先于”个人的学说的意义上。自然法论文与此相应，说伦理“如果不是个人的灵魂，就不能在个人身上表达出来。只有当它是一个普遍物和一族人民的纯粹精神时，它才是个人的灵魂。按照自然，肯定的东西先于否定的东西。或如亚里士多德所言，人民按照自然先于个人”。\footnote{《论自然法的科学探讨方式》，第396页。}这一基本原理的双重表达方式表明，黑格尔在进行其称作“伦理自然”的自然法构建时，同时以亚里士多德和斯宾诺莎为基础。\footnote{伊尔廷着重强调青年黑格尔的法和国家学说的这一重要的、迄今尚未得到足够重视的出发点，这是他的功劳。参见伊尔廷：《黑格尔对亚里士多德政治学的分析》（Hegels Auseinandersetzung mit der aristotelischen Politik），载《哲学年鉴》（Philos. Jahrbuch）第71卷(1963-1964)，第38-58页。}黑格尔把斯宾诺莎关于有限物的本质的一般教导\footnote{参见斯宾诺莎：《伦理学》，第一部分，命题八，附释一。黑格尔在《信仰与知识》第64页引用了这段话。}——有限物是否定物，它不能独自存在——与亚里士多德相对具体的也就是政治的主张——孤立的个人不是自足的，而是“在自然上”依赖于城邦——等同起来。我们在这里感兴趣的并不是那些在本体论上引发或者激发这种等同的理由。重要的是，黑格尔诉诸于斯宾诺莎，并不是考虑到他与亚里士多德迥然不同的自然法和国家法，而是只考虑到他的那种在某种程度上可以与亚里士多德的《政治学》或柏拉图的《理想国》相提并论的形而上学学说。是的，我们可以说，黑格尔之所以在批判自然法时利用斯宾诺莎的形而上学，主要是为了能够更为纯粹地阐明“绝对伦理”的结构，而其构造要素则由他借自古典政治学。\footnote{黑格尔自己并没有非常清晰地呈现奠定这种借用的基础的理念，谢林则在《学术研究方法讲座》（1802年）中对其进行了陈述：“就其内在生命而言，国家的科学建构并不会在后来的时代找到任何相应的历史要素。”《全集》第3卷，施罗特(M. Schröter)出版，慕尼黑，1927年，第335页；也见第315页。[谢林]这些地方[的说法]很可能已经考虑到了黑格尔的重构努力，或是参考了与黑格尔的谈话。同时也不排除黑格尔的论述也受到了谢林的影响。两人进一步的一致，比较谢林，同上，第314页（关于古代和现代国家之差异的原则），以及黑格尔，《论自然法的科学探讨方式》，第383页以下。}黑格尔同时也赋予[斯宾诺莎的]无限实体“按照自然先于”其分殊\footnote{参见《伦理学》，第一部分，命题，定义三和定义五。（中译文参见〔荷兰〕斯宾诺莎：《伦理学》，贺麟译，商务印书馆1983年版，第5、3页。——译者）}的学说以一种政治-实践意义。对斯宾诺莎来说，就像被理解为对“绝对伦理”的“理智直观”的“对上帝的理智的爱”的学说一样，无限实体学说并不具有政治实践意义。黑格尔明显影射斯宾诺莎道：“按照哲学关于世界和必然性的观点，万物都在上帝之中，没有任何个别性存在。对经验意识而言，由于每一个别的行动或思维或存在唯有在整体中才具有其本质和意义，因而这种观点便完全实现了。”\footnote{《伦理体系》，第465页，以及前一句：“因此在伦理中，个体以一种永恒的方式存在，其经验的存在和行动是完全普遍的存在和行动。因为不是个体，而是居于个体中的普遍的绝对精神在行动。”}黑格尔通过把斯宾诺莎吸收到对近代自然法的批判中，就在形而上学上升华了作为这种批判之尺度的古典政治学的那些基本原理。为黑格尔的自然法概念奠定了基础的“肯定的”自然概念，这个表达伦理整体绝对优先于其部分以及这种关系之必然性的概念，只会赋予个体一种完全“否定的”意义。\footnote{当在近代道德哲学和法哲学中表达出来的伦理意识在“普遍物和特殊物的同一”（其中普遍物为根据）之间“插入任何一个其他的个体性为根据”时，黑格尔就径直把它称为“非伦理的意识”。参见《伦理体系》，第466页；以及《论自然法的科学探讨方式》，第349页。}由此，黑格尔就在其法哲学的第一个发展阶段从古典政治学的立场出发对近代自然法进行了毁灭性的批判。黑格尔对斯宾诺莎形而上学的运用，使得这种批判彻底化了，更有甚者，它首先阻碍了黑格尔对近代自然法获得一种恰当的历史性理解。

\section*{2}

黑格尔批判各个自然法原理的起点和终点是，在这些自然法原理中，“个人的存在被设至首位和最高位置”。现代“主体的绝对性”在一种“比较低级的抽象”层次上出现在启蒙运动时期的幸福论当中；在“那些被称作反社会主义的”（霍布斯、卢梭）体系中，自然法把这种绝对性提升为理论。在康德和费希特的哲学中，这种主体的绝对性达到了对自身的理解，也就是对它的“概念”的理解。\footnote{《论自然法的科学探讨方式》，第343页以下。}在17世纪自然法中，这一原则把绝对伦理的环节歪曲为两种特殊的“本质性”，即自然状态和法权状态，在这两种状态中，个体性以不同的方式固定起来：在自然状态中，作为绝对自由；在法权状态，作为臣民的绝对服从。预设的自由以服从于一种对个人而言是外在的，并且就此而言其本身也是个别的、特殊的权力为条件。然而，个体性的取消不过是一种假象；实际上，它通过“服从的同一”关系不断地再生出来。对于结合起来的众多个体来说，伦理的结合(unio civilis)仍是一个异己之物，自然法“以社会和国家的名义”表象出来的一个无形式的、外在的统一。\footnote{同上书，第337页以下。}

黑格尔把这种“空名”与“绝对的伦理理念”进行比较。在绝对的伦理理念中，个别性并没有被固定起来，处于一种相对的同一性，[也不是]作为一种服从的关系（在这种关系中，个体性是某种完全被设定起来的东西），相反，“个体性自身纯属虚无，它与绝对的伦理威严纯为一物，——唯有这种真正的、生动的、并非服从的同一才是个人的真正伦理”。\footnote{《论自然法的科学探讨方式》，第338页。}就像在所有地方一样，在这里，个体也是否定物，其持存必须在绝对理念中“流动”起来。黑格尔很少像这一时期那样如此清楚地表达出古典城邦伦理的立场。他对[个人]与伦理整体的生动同一与近代自然法的国家理论的“服从的同一”进行了鲜明的对比，并把彻底否定自然法理论的出发点、“个人的虚无”作为前一种同一的前提。他在批判康德和费希特的自然法时又重复了这一点。对黑格尔而言，诸如自由、纯粹意志、人权这些概念都不过是“更为纯粹的否定物”\footnote{同上书，第340页以下。这种批判同时也是对他自己伯尔尼和法兰克福时期草稿所持立场的批判。青年黑格尔的这种立场，与卢梭、康德和费希特一道，想要让“从内心出发自我立法的不可转让的人权”抽象地生效。参见《早期神学著作集》，图宾根，1907年，第212页以下，另见第42、53、71、139页以下、173、188页以下、365页。参见《黑格尔来往书信集》(Briefe von und an Hegel)第1卷，J. 荷夫迈斯特出版，汉堡，1952年，第23页以下。对此立场的最早修正已经出现在《论符腾堡最近的内部情况》（1798年）和《德国宪制》（1800-1802年）中，然而这两个文本在我们这一研究的考察范围之外。}，它们的长处在于把最新伦理体系原则的“自为存在和个体性”更加清楚地表达了出来。在17世纪自然法中，“服从的同一”关系有时候前后不一，有时候又成为了简单的统治关系。\footnote{参见《论自然法的科学探讨方式》，第337页。}黑格尔以费希特为例表明，这种“服从的同一”关系会成为一种强制的体系，其全面性(Universität)使得“个别性与普遍性的同一”完全不可能。\footnote{同上书，第362页以下。}尽管观念论哲学的“伟大方面”在于，把权利和义务的本质与思维主体和意求主体的本质设定为同一，但是观念论哲学并不忠实于这种同一性原则。不仅这个原则的有效性由于限于行为的主观道德而被打破，分裂也在合法性中同样被绝对地设定，而且观念论哲学认为[道德和合法性]两者具有同等价值，并且没有丝毫关系。主体与权利和义务的本质的同一的可能性（即道德性）与不同一的可能性（即合法性）相对立，自由的体系与强制的体系相对立。\footnote{参见《论自然法的科学探讨方式》，第360页以下。}

因此，尽管黑格尔把道德的思想尊为康德和费希特哲学的“伟大方面”，但是对他而言，道德只不过以更为纯粹的形式表现了个体性的否定物。因为作为单纯的可能性来经验的并且与合法性分离开来的意识和义务的同一把道德主体抛回到其个体性中。诚如黑格尔所言，构成了道德和合法性两者之本质的，以及两者在其中是一个东西——“伦理自然”——的，并不是“肯定的绝对物”，而是以前述的形式与之完全对立的否定的绝对物，或者说“绝对概念”。现在，黑格尔同时也把观念论自然法体系内部的这种对立以及道德原则中对否定物的固定与“绝对伦理”的立场进行对照。在方法论上，这是通过诉诸“理智直观”而发生的。在黑格尔看来，在理智直观中并不存在可能性和现实性、概念和存在（“自然”）的分离。因为就像古代城邦理论一样，这种历史性地在其诸环节中重建了古代城邦理论的伦理自然的直观，包含了现实东西的当下性，一个作为“活的关系和绝对当下”的“这一个”。\footnote{参见上书，第357、359页。此外，这里我们还可以注意，黑格尔对从谢林那里接受过来的合乎历史地占有古典伦理学和政治学的方法进行了多大程度的创造性发挥。}在这里，黑格尔为了对他把古典政治学和现代自然法进行对比的方法作辩护，转向了语言：“表示‘伦理’的希腊词，还有德语词，都卓越地”暗示了语言的一般预设，即在“绝对伦理的自然”中存在着一种“普遍物或伦常”，“而最近的伦理体系由于以一种自为存在和个别性为原则，就无法做到使用这些词语而不言明它们的联系；这种内在的暗示如此有力地证明，为了标明其本质，这些最近的伦理体系就不能滥用这些语词，而是采用了道德一词[……]”。\footnote{《论自然法的科学探讨方式》，第396页。}将伦理的理智直观与伦常的“这一个”的联系起来，其结果就是必然向古典自然法靠拢。

这种靠拢首先表现在黑格尔对道德学的评价上。尽管绝对伦理作为德行论直接就是个人的伦理，然而对黑格尔来说，个体乃是完全的否定物，[由此]绝对伦理就以否定的形式显现在个体中；不仅现代主体性的反思性道德，而且渗入伦理整体的联系之中，并且表现于外的个体的德行（黑格尔想要按照道德学的古典概念模式将这种个体的德行理解为对道德学所做的自然描述\footnote{参见上书，第399页。}），也是“否定性的伦理”。道德学与自然法（在一种构建“伦理自然”法意义上的政治学）的形式差别就在于此；我们不能像在康德和费希特那里那样这样理解这种差别：“仿佛它们是分离的，道德学被自然法排除掉了，而是它们的内容完全就在自然法之中。”\footnote{参见《伦理体系》，第465页；另见《论自然法的科学探讨方式》，第396页。}正如自身否定的东西的领域、作为普遍物的单纯可能性的个体的伦理归于道德学之下一样，归于自然法下面的是“真正的肯定物”的领域、它的现实性。在这一时期，黑格尔把这种现实性叫作“实在的绝对伦理”或者“伦理自然”。在它的由一族人民设定起来的法中构建这种“自然”，就是黑格尔为自然法规定的任务。如果像在近代[自然法的]科学探讨方式那里那样，当“自然法的规定由”否定的东西、道德的抽象物、“纯粹意志以及个体意志的抽象物以及诸如强制、通过普遍的自由概念对个人进行限制等抽象物的综合表达出来时”，那么自然法的概念就只能是一个否定的概念，“一种自然的非法[Naturunrecht]，因为当这些作为实在性的否定物为[自然法]奠基时，伦理的自然即已处于彻底的败坏和不幸之中了”。\footnote{《论自然法的科学探讨方式》，第397页。[引文中的“自然的非法”(Naturunrecht)，德文原文误作“自然法”。兹据黑格尔原文（见G. W. F. Hegel. Gesammelte Werke. Band 4. Felix Meiner Verlag Hamburg 1968, p. 468）改正。——译者]}

\section*{3}

对这种自然法概念以及与之相联的将伦理设定在一族“人民”（对人民的理智直观只不过是非常明显地模仿了雅典城邦的轮廓）中的方法论立场，黑格尔并未坚持多久；他在自然法论文发表之后不久就修正了自然法概念，最后在耶拿末期完全放弃了它。黑格尔学界只是偶尔注意到这种概念上的转变。\footnote{参见F. 达姆施泰特(Darmstaeter)：《作为社会权力的自然法与黑格尔的法哲学》(Das Naturrecht als soziale Macht und die Rechtsphilosophie Hegels)，载《智慧》(Sophia)第5卷(1937年)，第215页，注释13。达姆施泰特只是发现黑格尔主张过各种不同的自然状态观，但没有对这种不同进行追问。海林也注意到黑格尔“在非常不同的意义上”使用“自然”的概念，一种无助的状态（海林：《黑格尔》第2卷，第392页以下）。认为黑格尔后来在概念的使用上仍然摇摆不定（第393页），则是错误的。}这种引人注目的概念转变与黑格尔在1803-1804年与1805-1806年的耶拿讲座之间完成的思想基础上的一般转变有关。\footnote{对此，参见罗森克兰茨：《黑格尔传》，第178页以下；F. 罗森茨威格：《黑格尔与国家》第1卷，慕尼黑/柏林，1920年，第183页。}在这几年间，黑格尔似乎重新研究了费希特哲学，同时从谢林的术语和方法中摆脱了出来。这也意味着同时放弃迄今为止以亚里士多德和斯宾诺莎为准的自然法构想。就在细节上对黑格尔体系和概念形成上的这种转变的前提和结论进行考察而言，那就要一方面在与自然、“概念”和法的联系中，另一方面在它们与“个体性”以及“个体性”在体系中的价值关系中去寻求。黑格尔之所以在伦理体系草稿和自然法论文中把个体的存在规定为完全的否定物，并且沉没在整体的“伦理自然”之中，是因为他在这一时期仍然主要把否定物等同于虚无，等同于趋于纯粹的毁灭。这样，为近代自然法奠定其“社会和国家”演绎之基础的个体的存在，就在第一个体系草稿中出现在“否定物或自由或犯罪”这一标题下面。\footnote{参见《伦理体系》，第446-460页。}它的主题是：物理的毁灭、抢劫、偷盗、征服、谋杀、报仇、斗争、战争。体系的这一部分位于(1)“自然伦理”（按照关系的绝对伦理）和(3)（绝对）伦理之间，而不触及两者中的任何一个。倘若伦理作为“肯定物”直接设定起来，那就必须事先拒绝一种过渡的可能性。在后来的那些体系草稿中，这一部分获得了一种完全不同的功能。在1803-1804年讲座中，这一部分尚处于自然伦理（家庭）和人民之间，但是已经失去了其自己的体系地位。\footnote{参见《耶拿实在哲学》第1卷，J. 荷夫迈斯特出版，莱比锡，1932年，第226-232页。这一章没有自己的标题。对此，参见伊尔廷（本篇注释21）非常中肯的观察，第56页。}在第一个体系草稿和自然法论文中，由于“人民”的自然设定而被排除在外的那种向绝对伦理的过渡，已经在这里出现。最后，在1805-1806年讲座中，体系的这一部分压缩为短短几页，而且还是出现在完全不同的位置，或者说页边。\footnote{《耶拿实在哲学》第2卷，莱比锡，1931年，第210-212、219-221页（即荷夫迈斯特对第237页注释3的页边补充）。}

然而，黑格尔在自然法论文中就已经暗示了从个体的否定性存在过渡到伦理的肯定性存在的最初可能性。个体本身所是的否定物，可以由个体所设定，对其各种规定性的否定也可以被绝对地否定。这种从一切限制中抽身出来的能力，尽管也是一种否定物，但却是一种这样的否定物，通过它，个体存在成为“被绝对吸收到概念之中的个体性、否定的绝对无限性、纯粹自由”。这样理解的自由、个体的“概念”，作为否定的绝对物，是绝对物自身的环节。也就是说，包含在绝对伦理中。绝对伦理出现在“自由人等级”。自由人等级在生死斗争的活动中证明了纯粹的自由。但一个肯定物作为其条件要先于这种否定的活动，[这就是]“绝对的伦理事物，亦即属于一族人民之物[……]个人与它的同一唯有在否定物之中、通过死亡的危险才能得到确凿无疑的证明”。\footnote{参见《论自然法的科学探讨方式》，第370页以下；《伦理体系》，第465页以下。}这里表现出来的局限是黑格尔当时的立场；过渡的可能性只保留给了自由人，而非自由人等级，由于其劳动并不合乎于“概念”\footnote{参见《伦理体系》，第473页：“他们并不在无限的概念之中，通过无限的概念，这个对他们的意识而言不过是作为外在物而被设定之物，就会完全是他们绝对的、自身的、推动他们的精神，这个精神克服了他们的一切规定性。他们的伦理自然达到了这种直观，这种好处是第一等级提供给他们的。”}，则囿于个体性之中。此外，第一等级的否定性活动不能从自身中产生出任何恒常之物（“肯定物”）；相反，伦理的肯定的持存总是已经预设好了的。

黑格尔在1803-1804年讲座中还在做出这种预设。在实现从个体性到普遍性过渡的地方，讲座还几乎与《伦理体系》和自然法论文逐字对应：“个人作为人民的成员是一个伦理的存在，他的本质是普遍伦理的生动实体，他作为个别的东西，仅仅作为一种被扬弃了的东西，[是]一种观念性形式的存在者。伦理在其生动的多样性之中的存在就是人民的伦常。”\footnote{《耶拿实在哲学》第1卷，第232页。}这仍然完全是古典城邦伦理的立场，以及前述所有方法论的和体系与内容的结果。很明显，只要坚持这种立场，与近代自然法理论的关系就会一直都是否定性的。黑格尔在边注中写道：“没有和解，没有契约，没有任何默示或明示的原始契约。个人[不仅必须]放弃他的自由的一部分，而且必须完全放弃[自身]。”\footnote{同上。}尽管如此，过渡不再局限于“自由人等级”。否定性的活动不再是为了保存伦理整体而斗争的活动，而是个体为了相互承认的斗争。\footnote{在《伦理体系》，这个主题尽管没有特定的价值，但已经在发挥作用了：一方面，它出现在论述交换、财产、货币和贸易的承认领域（即个体性成为普遍）之后，另一方面它位于家庭关系之前。在家庭关系的“无差异”中，为了承认的斗争的结果（主奴关系）得到了扬弃。参见《伦理体系》，第440-443页；另见第495页以下，在这里，法权体系事后被解释为承认的形式体系（第497页）。}只有到了现在，那个在体系草稿中只占据了一个漠不相干的位置，并且在自然法论文的过渡理论中束缚在一个等级的偶然性上面的“否定物”，方才获得了那种构成黑格尔法思想最为重要的基础之一的功能：个别性与普遍性之间的中介。在这里，为了承认的斗争首先构成了家庭的“自然”伦理与人民的“绝对”伦理之间的中介，同时，人民的伦理一如既往地来自于对亚里士多德和斯宾诺莎的范畴的改写。\footnote{参见《耶拿实在哲学》第1卷，第232页：“一族人民的绝对精神是绝对的普遍要素，是将所有单个的自我意识交织于自身的以太，是绝对简单的、活生生的、唯一的实体。”}这样，个体得到承认，也就包含了它被接纳到绝对的伦理实体之中，[从而]意味着个体自身的消失\footnote{同上：“因此，这种绝对意识（在这个阶段，是为了承认的斗争的结果——笔者）是作为个体的意识的被扬弃了的存在[……]。它是普遍的、持存着的意识；它不是无实体的个体性的单纯形式，而是这些个体不再存在；它是绝对实体。”}。

尽管如此，却由此迈出了克服古代城邦理论的抽象方法论立场以及重新评价现代自然法的第一步。这一点可以从1805-1806年讲座的相关段落中很好地看出来。在这些段落中，黑格尔从体系、方法和内容上得出了这种出发点的结论。谢林的方法及其术语的含义完全消失不见了；黑格尔现在使用的范畴不再是“肯定物”和“否定物”、“伦理自然”及其法、“直观”和“概念”的对立等，而是起源于“自我”的“理智”和“意志”。1803-1804年讲座中隐现在“承认”的中介功能中的东西，现在变得完全清楚了：黑格尔在耶拿时期后半段重新对费希特进行了研究，并大约从1804年开始在他的精神哲学和自然法的讲座中回到了费希特。\footnote{对费希特“主体性形而上学”的积极继受（这种继受规定了黑格尔耶拿讲座一些章节的结构和安排）肯定在耶拿逻辑中发挥了重要作用。在出版者（H. 埃伦贝格、G. 拉松）看来，耶拿逻辑属于黑格尔最早的、时间定为1801-1802年的讲座草稿。然而，其理论基础“自我”作为个别性和普遍性的综合的观点（参见《耶拿逻辑学、形而上学和自然哲学》，G. 拉松出版，莱比锡，1923年，第161、163页以下、178页以下），在1801到1803年间撰写的论文中以及部分地在1803-1804年的讲座中都还绝对没有达到，这种情况在前面探讨过的黑格尔自然法理论第一个文本（即《论自然法的科学探讨方式》）中已经表述出来了。只有到了1805-1806年讲座才出现了转向。这一讲座的安排也与《耶拿逻辑学》的相关章节（第三章 主体性形而上学，第161页以下）明显一致。}现在，其探讨方式暴露了费希特的直接影响的承认主题，\footnote{比较费希特：《自然法权基础》（1796年），《费希特全集》第3卷，第44页以下、85页以下，与黑格尔：《耶拿实在哲学》第1卷，第226页以下；《耶拿实在哲学》第2卷，第194页以下、205页以下。对于黑格尔而言，费希特对“承认”作为（“现实的”）行动与“概念”的综合解释提供了一种结合的可能性，在这种综合中，各个个体彼此交互地“设定”起来。}被明确地引入近代自然法的提问方式中：“这就是习惯称之为自然状态的关系：各个个体之间的自由、平等的存在，自然法应该回答，按照这种关系，个体彼此间具有怎样的权利和义务，对他们[作为]按照他们的概念而且是独立的自我意识而言，他们的行为的必然性是什么。”\footnote{《耶拿实在哲学》第2卷，第205页。}黑格尔自己用众所周知来自霍布斯的论点：必须走出自然状态(exeundum e statu naturae)，进行回答。他已经在耶拿教授资格论文的第九个命题中运用过这一论点，尽管论证意图有所不同。\footnote{参见黑格尔：《最早出版作品集》(Erste Druckschriften)，G. 拉松出版，莱比锡，1928年，第404页：“自然状态是不义的，因此必须走出这种状态。”(Status naturae non est injustus, et ob eam causam ex illo exeundum.) 众所周知，霍布斯和斯宾诺莎恰恰论证的是相反的观点：应当离开自然状态，因为它是不义的。黑格尔在注释53中所提到的地方赞同这一说法。}康德和费希特也持有相同的论点，即个体“在自然上”不具有任何权利和义务，只有法权状态、“国家”才要求权利和义务具有现实性（效力）。

[在他们这里]“设定起来的”，不再是一种无论怎样理解的“自然”，而是“概念”，正是这种概念应当达到它的“法”：“彼此自由的自我意识概念被设定了起来，但也恰好只是概念而已；概念，正因为它是概念，毋宁必须实现它自身，也就是说，扬弃与其实在性相对立的、处于概念的形式中的自我。”\footnote{《耶拿实在哲学》第2卷，第205页。之前的黑格尔学界想要用在许多方面都不充分的主要与《精神现象学》发展阶段相连的公式“伦理的去浪漫化”去把握我们所描述的这种黑格尔自然法构想的转向。参见海姆(R. Haym)：《黑格尔及其时代》(Hegel und seine Zeit)，柏林，1857年，第207页以下；梅茨格(W. Metzger)：《德国观念论伦理学中的社会、法和国家》(Gesellschaft, Recht und Staat in der Ethik des deutschen Idealismus)，海德堡，1917年，第310页以下。}自身实现的概念无非就是承认的运动；这个运动没有直接过渡到在1803-1804年讲座中依然还是唯一的肯定物的伦理的“绝对实体”，而是过渡到“一般伦理”，并且是过渡到它的直接性，也就是法。而伦理观上的这种值得注意的转变意味着对黑格尔迄今为止的自然法构想基础的颠倒。承认的运动得以产生的“概念”并不是那种将字面意义上的“无”从自身中释放出来的[作为]个人的“纯粹自由”的否定物；因为概念的运动自身产生出了肯定物的要素，在其中个体性达到了一种恒常物，也就是法：“法是人格(Person)彼此对待的关系，是人格的自由存在的普遍要素或者说规定，[也就是]对它的空洞的自由的限制；我不必为我自己策划和产生这种关系或限制；相反，对象自身就是一般的法，亦即得到承认的关系的这种产生[……]被承认者通过他的存在而被人承认为直接有效，然而，这种存在也正好是从概念产生出来的。”\footnote{《耶拿实在哲学》第2卷，第206、212页。}概念的运动打破了自然、伦理和法迄今为止的（“直接的”）联系，并且与自然相对，在法中将自身固定下来。黑格尔继续展开的与自然法的自然概念的争议也不能掩盖这一事实，即这场争执现在已经取得了一种新的标准。这个标准本身归功于观念论晚期的自然法，并且恰恰与迄今为止的自然法相对立。这一标准就是普遍的法权能力，作为纯粹概念（即人的自我）的“人格”。“人格”的发现以与一切既定的自然秩序及其“规律”决裂为前提。“法包含了纯粹人格，纯粹的被承认存在。因而人格并不在自然状态，而是沉入定在，由此他是一个在其概念中的人；但在自然状态中，他并不在他的概念中，而是作为自然存在，在其定在中。问题直接地自相矛盾——我考察的是在其概念中的人，而不是在自然状态中的人。”\footnote{《耶拿实在哲学》第2卷，第205页（黑格尔的边注）。荷夫迈斯特并没有完全正确地读懂第二句；“[……]沉入定在”后面似乎应当打逗号，因此读作“借此，他是人，[他是]在他的概念中”。}

向费希特靠拢首先表现在我们这里感兴趣的讲座段落的体系安排上。在1803-1804年讲座中，承认的运动还是意识的一个“潜能阶次”，它在伦理的“绝对意识”中扬弃其自身。而在1805-1806年讲座中，意识的基础则是自身设定为“理智”和“意志”的“自我”的自发性。这种自发性就是上面所言的概念的运动，承认表达了这种运动的“肯定的”形式。黑格尔在这一章把“自我”等同于“概念”，实非偶然：“[……]自我，概念自身[就是]运动。自我通过他者而运动，将他者转化为自身运动，并且反过来把这种自身运动转变为他在，自我的对象化，以及自为存在。”\footnote{同上书，第204页（边注）。}黑格尔在耶拿《逻辑学》中已经把自我理解为个别性和普遍性的统一，并且与承认范畴联系起来。\footnote{参见《耶拿逻辑学》，第164页以下，第171页。}这里，自我不再显现为居于自然的“生动”统一之下的空洞的统一，也不显现为对立环节的单纯统一，而是显现为这些对立环节的一种通过进行设定的活动创立起来的中介，概念的自身中介。\footnote{对此，参见施瓦茨(J. Schwarz)：《黑格尔的哲学发展》(Hegels philosophische Entwicklung)，美茵河畔法兰克福，1938年，第339页以下。施瓦茨指出，要对《伦理体系》第431页、《信仰与知识》第56页以下与《耶拿实在哲学》第2卷第190页和195页注释5进行对照。}因此，黑格尔最初赋予传统“自然法”之名的积极意义，必须被抛弃。与这一出发点相应，黑格尔在耶拿《逻辑学》中已经把精神抬高到自然之上，并在自然法论文中把“物理自然”和“伦理自然”区分开来。\footnote{参见《耶拿逻辑学》，第192页以下；《自然法》，第395页。}但在涉及自然法的奠基和概念时，他却并没有利用这种区分，其原因要在片面的古典政治学和斯宾诺莎的实体形而上学的定位上面去找；由此，黑格尔就为自然法赢得了“伦理自然”的“法”的意义，并将其等同于一族“人民”的实体性伦理。\footnote{《论自然法的科学探讨方式》，第346页以下。}直到在1805-1806年讲座，在对费希特和——有很大可能——卢梭重新进行研究之后，\footnote{参见《耶拿实在哲学》第2卷，第213页以下、244页以下，这里尽管没有出现卢梭的名字，但是明显涉及《社会契约论》的表述。在之前的草稿中，“公意”的概念完全没有起作用。}黑格尔才超出了这种立场，并与一种以现代理论的“否定性的”自然概念为基础的自然法概念达成一致。他构建的自然法并不是一种无规定的“伦理自然”的法，而是承认运动的法。在这种承认运动中，个体性的持存、“法”得以产生出来。这种运动的必然性在于个体的“概念”，正如黑格尔所言，这种运动“是概念自己的运动，而不是与内容相对立的我们思维的运动。作为承认，概念自身就是运动，而这种运动正好扬弃了概念的自然状态：概念是承认；只有自然的东西存在，它并不是精神的东西”。\footnote{《耶拿实在哲学》第2卷，第206页。这个地方必须与第242页对照：“精神是众多个体的自然，是他们直接的实体，以及这一实体的运动和必然性。”}进行承认的运动——它的“概念”构成法——在个体性的方面把一种“意志”作为结果。这种意志是“认识着的”意志，或是一个“普遍物”，——卢梭的“公意”，它的直接现实性就是具有法权能力的“人格”。\footnote{《耶拿实在哲学》第2卷，第212页：“现在，这个认识着的意志是普遍物。它是得到承认的存在。它在普遍性的形式下与自身对立，它就是存在、一般现实性，而个体、主体是人格。个体的意志就是普遍的意志，普遍的意志就是个体的意志，一般的伦理，而它直接就是法。”相反的观点，参见拉伦茨(K. Larenz)：《德国观念论的法哲学和国家哲学》(Die Rechts- und Staatsphilosophie des deutschen Idealismus)，见《哲学手册》(Handbuch der Philosophie)，第四部分，慕尼黑和柏林，1934年，第153页。他认为意志的出现始于《精神现象学》。}把这两个范畴纳入伦理体系，表明黑格尔与古典政治学构想及其自然概念的决裂，而在1802-1804年间，他是把它们同现代自然法抽象地对立起来的。随着对作为伦理事物的直接性的“法”的承认以及将法纳入“概念”之中，个体性的环节，也就是现代自然法的出发点，便得到了辩护。由此，在经历了耶拿前期的动摇之后，黑格尔最终重新找回了他曾在1790年代追随过的卢梭、康德和费希特的自然法立场。\footnote{参见《书信集》第1卷，第18、23页以下；《早期神学著作集》，第42、53、62、72、139页以下、173、188页以下、365页。黑格尔于1798年研习了康德的法权论，1796年研读了费希特的自然法。参见罗森克兰茨：《黑格尔传》，第48、87页；梅茨格：《德国观念论伦理学中的社会、法和国家》，第331页以下。}黑格尔在耶拿末期获得了这样的洞见：“一切”都取决于“不仅把真实的东西理解和表述为实体，而且同样理解和表述为主体”。\footnote{参见黑格尔：《精神现象学》，J. 荷夫迈斯特出版，汉堡，1952年，第19页。（中译文引自《精神现象学》上卷，第10页。——译者）}通过“概念”（即自我）——在概念中，个别性和普遍性得到了中介——对个体的存在所做的这一辩护，就是这种洞见的结果。这尤其适用于自然法的观点；自然法的真理是（再次引用《精神现象学》序言，黑格尔在其中保证了他的哲学的新方法论视域）：只有当主体“是设定自身的运动时，或者说，只有当它是自身转化与其自己之间的中介时”，它才是“现实的”。\footnote{《精神现象学》，第20页。（中译文引自《精神现象学》上卷，第11页。译文稍有改动。——译者）}自然法不再以“伦理自然”和黑格尔最初将之等同于希腊城邦的实体性伦理的“伦理自然”的法，而是以在其运动中渗透了伦理实体的法的“概念”为对象。

\section*{4}

设定自我的概念打破了自身“自然的”直接性和实体性伦理的“自然”。黑格尔只有在洞见了这种自我设定的概念的运动之后，才作为完成了的“伦理学家”出现在我们面前。这种双重运动相应于这样一种基本立场，在这种立场中，黑格尔第一个自然法构想的困境及其克服历史性地发展出来。现在，在耶拿末期，由于黑格尔也从自然法的立场出发批判古典政治学，从古典政治学的立场出发来对现代自然法进行批判便颠倒了过来。这第一次明确地出现在1805-1806年讲座中。一方面，黑格尔把对从卢梭那里接受过来的从个体意志中构建起来的公意理论的解释追溯到亚里士多德的“整体按照自然先于部分”的思想，由此得出，对于个体意志而言，公意\footnote{即普遍意志。——译者}乃是“第一位的东西和本质”。\footnote{参见《耶拿实在哲学》第2卷，第244页以下。}另一方面，尤其是由于吸收了那种自然法理论，黑格尔又把分裂、作为“近代的更高原则”的“个体的自身绝对认识”与古典城邦伦理中的“普遍物与个体的直接统一”对立起来。\footnote{《耶拿实在哲学》第2卷，第251页。}其后，黑格尔即用这个产生了“概念”双重运动的标准去衡量[古典]政治学和自然法，并由此理解到两者的局限。在哲学史讲座中，我们又遇到同一个标准。这种一致绝非偶然。因为众所周知，黑格尔在1805-1806年第一次讲授了这门课程，并且后来也利用了耶拿手稿。在关于柏拉图的讲座中，奠定国家(Politeia)的基础的“实体性的东西、普遍的东西”与“就自然而言的法，[……]在个体上面与为了个体的法”形成了对照。\footnote{黑格尔：《哲学史讲演录》，《全集》第14卷，柏林，1832之后，第269页以下。这里“伦理自然”这一术语还在使用，但显然意思已经转变为“在其合理性中的自由意志”（第269页）。（中译文参见〔德〕黑格尔：《哲学史讲演录》第二卷，贺麟、王太庆译，商务印书馆1960年版，第243页以下。——译者）}黑格尔把这种近代自然法称为一种“对实在的实践本质、对法的琐碎抽象”\footnote{参见《全集》第14卷，第290页以下。}；而另一方面，他又强调柏拉图城邦理念的局限，其对个体性、个体意识、直到近代自然法才使之有效的“人格”\footnote{同上书，第295页。（中译文引自《哲学史讲演录》第二卷，第266页。译文略有改动。——译者）在对亚里士多德《政治学》的讲座（第400页）中，再次重复了黑格尔批判的双重方向。比较对柏拉图《理想国》的评断（第275页）与《耶拿实在哲学》第2卷，第251页：“柏拉图并没有提出一种理想，而是内在地把握到了其时代的国家。但是这种国家没落了[……]，因为它剥夺了绝对个体性的原则。”（边注）}的排斥：“与柏拉图的原则正相反对的是个人的自觉的自由意志原则，这一原则近来特别被卢梭提到很高的地位：认为个人本身的意志、个人的表现是必然的。”\footnote{参见《全集》第14卷，第295页。}

对古典政治学和现代自然法的局限的认识，同时也是以一种对近代自然法本身的特定历史理解为前提的，因为黑格尔的批判标准并不是从随便哪个作者那里取来的，而是取自作为批判的背景的卢梭。《哲学史讲演录》和《历史哲学讲演录》都表明，在黑格尔这里实际上存在着这样的理解。尽管黑格尔也认同18和19世纪广为传播的意见，即自然法的“启蒙运动”是从格劳秀斯的《战争与和平法》（1625年）开始的，但同时黑格尔也指出，支配格劳秀斯著作的仍然是斯多亚—西塞罗的自然目的论。\footnote{参见《哲学史导论》(Einleitung in die Geschichte der Philosophie)，第3版，J. 荷夫迈斯特出版，汉堡，1959年，第159页以下；《哲学史讲演录》（《全集》第15卷），第439页以下、445页以下、502页；《世界历史哲学讲演录》，拉松出版，莱比锡，1920年，第917页以下。黑格尔指出的古老自然法的“naturaliter insita”，显然指的是西塞罗：《论法律》(Legg.)第1卷，6.18。这里，自然法(lex naturalis)作为“根植于自然的、指挥应然行为并禁止相反行为的最高理性”(“ratio summa insita in natura, quae jubet ea, quae facienda sunt, prohibetque contraria”——中译文引自〔古罗马〕西塞罗：《国家篇 法律篇》，沈叔平、苏力译，商务印书馆1999年版，第158页。——译者)出现。也参见《论法律》第1卷，12.33。（中译文参见同书，第165页。——译者）}按照这种自然目的论，权利意识对人而言是“自然的”或者说与生俱来的(“natura insitum”)，因此，无论是格劳秀斯、普芬道夫和沃尔夫，还是苏格兰学派的道德哲学，都未曾与法传统决裂。而在黑格尔看来，摧毁了格劳秀斯仍然接受的“自然”与法之间的奠基关系的真正“革命性的”自然法理论来自霍布斯的著作，因为他最先“试图把维系国家统一的力量、国家权力的自然回溯到内在于我们本身之中的原则，亦即我们承认是我们自己所有的原则”。\footnote{《全集》第15卷，第442页。（中译文引自《哲学史讲演录》第四卷，第157页。译文略有改动。——译者）}在霍布斯那里，法和社会的理论设定一种自然秩序(ordo naturalis)的条件是，这种自然秩序需要规定使得离开它成为必然的那些条件。由于霍布斯，自然法的传统名称中产生了“歧义”。对黑格尔来说，这种歧义至关重要，并在解释《论公民》时注意到了这种歧义：“‘自然’这一名词具有歧义，一方面，人的自然指他的精神性、合理性而言；另一方面，他的自然状态则指人按照他的自然性而行动的状态而言。”\footnote{《全集》第15卷，第443页。（中译文引自《哲学史讲演录》第四卷，第159页。译文略有改动。——译者）}在黑格尔看来，卢梭、康德和费希特推进了这种歧义，加剧了霍布斯与传统自然法理论的决裂。卢梭把人的“精神性、合理性”理解为人的自由，自由使人从自然中走出来，被视为人与动物区分开来的绝对标志。\footnote{同上书，第527页。（中译文引自《哲学史讲演录》第四卷，第233页。——译者）黑格尔引用了《社会契约论》第1卷第四章的著名段落；参见第一章、第八章、第十一章；第3卷，第九章，注释；《论人类不平等的起源》，K. 维甘德出版，汉堡，1955年，第106页以下。此外，参见《世界历史哲学讲演录》，第920页以下。}如黑格尔所言，我们必须认为卢梭的原则：人是自由的，并且建立在“公意”基础上的国家乃是这种自由的实现，“是正确的”。当卢梭从个体意志、个体自然性的自由倾向“共同设定”公意时，他的“歧义”方才产生。由此，“自然的自由”又取消了卢梭在公意概念中所思考的“作为完全绝对物的自由”。\footnote{参见《全集》第15卷，第528页：“这些原则表达得很抽象，我们必须把它们正确地找出来；可是歧义立即就发生了。人是自由的，这当然是人的实质本性；这种本性在国家里不但没有被扬弃，事实上倒是开始被建立起来了。本性的自由，自由的秉赋并不是现实的；因为国家才是自由的实现。”（中译文摘自《哲学史讲演录》第四卷，第234页。——译者）另参见《哲学全书》，第163节，补充一（《全集》第6卷，第360页）；《法哲学》，第258节。}尽管如此，同时卢梭也使人意识到了，自由应是“人的概念”。\footnote{参见《全集》第15卷，第528页。}即对黑格尔来说，从中可以得出“向康德哲学的过渡”。康德哲学通过把自然的立法与自由的立法、经验的意志与“自由的和纯粹的意志”彻底分开，终结了自然法概念中的“歧义”。\footnote{参见《全集》第15卷，第529、552页以下、590页。}

在先验哲学的立场上，人的“概念”才达到了对其自身的理解。“自我意识的单纯统一”、“自我”乃是从一切既定的自然秩序中摆脱出来的、“不可穿透的、完全独立的自由和一切普遍的规定也就是思维规定的源泉”。\footnote{《世界历史哲学讲演录》，第922页。}“精神的东西的意识”成为了自然法理论的基础；自然法理论发现了一种“国家的思想原则[……]这种原则不再是诸如社会性冲动、财产安全的需要之类的任何一种意见的原则，也不是像君权神授这样的虔诚原则，而是与我的自我意识同一的确定性原则”。\footnote{同上书，第924页。}

要理解黑格尔《法哲学原理》的形而上学基础及其在终结自然法中独特的模糊立场，我们就必须要知道他对近代自然法的这种历史评价。这种评价在根本观点上偏离了自然法论文。就像现在自身表明的那样，自1805-1806年起就构成了黑格尔法思想的视域的“概念”的辩证运动，绝不是抽象的、空无内容的公式；相反，这种运动的内容就是意志的绝对“自由”。自从卢梭、康德和费希特以来，这种自由就被提升为一切法的原则。卢梭早就认识到，这种法原则隐含了自然法理论的标准乃是对自然的否决。他区分了两种自然法：(1)“本来意义上的自然法”(droit naturel proprement dit)，这种自然法以一种自然情感为基础（同情的自然法），并使自然状态中的人能够组成社会；(2)经推理而来的自然法(droit naturel raisonné)，这种自然法在国家建立之后，在市民状态中起支配作用。\footnote{参见《社会契约论第一版》(Première Version du Contract Social)第1卷，伏汉(C. E. Vaughan)编，曼彻斯特，1918年，第493页。}卢梭拒绝一种“理性的”即以自然秩序为范本的自然法，也就是格劳秀斯及其18世纪的后继者们都从中获得国家法上效力的那种自然法。卢梭和康德、费希特在这一点上是一致的，尽管康德、费希特还比卢梭更进了一步。康德偶尔也会指出黑格尔在霍布斯讲演中提到的这种自然概念的“歧义”。康德在其“道德哲学反思录”中写道：“没有市民秩序，全部自然法就不过纯然是一种德行论，并且只是作为一种可能的外在强制法律的计划，而具有一种法的名称[……]因为‘自然法’这个词曾如此歧义地使用，因此我们必须小心避免这种歧义。我们要区分自然法(Naturrecht)和自然性的法(natürliche Recht)。”\footnote{反思录，第7084条（约1776-1777年），康德，《全集》（科学院版）第19卷，第245页。}自然法是内在于市民秩序的理性法，而自然性的法则是建立在自然的个别规定上面的法，依照其内容要么是私法要么是公法。不同于卢梭和康德，费希特不仅把这个传统的名称劈分为两个方面，而且完全抛弃了这个名称。因此，费希特毫不怀疑“除了在一个共同体中，除了在实定法之下，根本就不可能存在人们经常运用的意义上的自然法”。\footnote{费希特，《自然法权基础》（1796年），见《全集》第3卷，第148页；参见第55、92页以下。《遗著集》(Nachgelassene Werke)第2卷，第498页以下。}

在其法哲学思想的第三个阶段，黑格尔从这种特定的现代自然法问题中引出了概念历史上必然形成的结论。他在多大程度上向霍布斯、卢梭、康德和费希特的立场靠拢，或许最明显不过地表现在，他在这一时期赋予前述自然和法的关系上的“歧义”以何种意义。他在1817年的《哲学全书》中写道：“迄今为止，自然法这一表达，对哲学法学来说已经成为了一种习惯表达。这一表达具有歧义：法是不是规定为似乎由直接的自然植入的，还是指的是，它是由事物的本性，也就是说由概念来规定自身的。第一种意义是之前习惯所指的意义，以至于同时还虚构出一种自然法在其中应当有效的自然状态。”\footnote{黑格尔：《哲学全书》，海德堡，1817年，第415节。（中译文参见〔德〕黑格尔：《哲学科学百科全书》，薛华译，上海世纪出版集团2002年版，第301页。——译者）对此，参见福莱希曼(E. Fleischmann)：《黑格尔的政治哲学》(La philosophie politique de Hegel)，巴黎，1964年，第17、113页以下。}通过法学术语“事物的[自然]本性”的中介作用，“概念”完全把“自然”的意义集中于自身。这里，概念的含义与耶拿末期是一样的：法，就其把自身及其一切规定都建立在“自由的人格性”之上而言，乃是“一种自我规定，这种自我规定毋宁是自然规定的对立面”。\footnote{《哲学全书》（1817年），第415节。（中译文参见《哲学科学百科全书》，第301页。——译者）也见黑格尔：《纽伦堡著作集》，J. 荷夫迈斯特出版，莱比锡，1938年，第156、170页以下。特别是第284页[《哲学百科全书》(Philosophische Enzyklopädie)，第181节]：“精神作为自由的、具有自我意识的本质，乃是自身等同的自我，在其绝对否定性关联中首先进行排斥的自我，个体的自由的本质或人格。”}现在，在黑格尔最终理解了近代自然法的历史意义之后，他就比耶拿时期更为明显地把概念的抽象自我运动等同于摆脱了所有“自然性的”规定的“人格”的自我决定。自由，观念论自然法所阐发的人的无限的、不为任何外物所限制的概念，乃是法的唯一原则。

黑格尔在柏林时期的法哲学讲座中反复讲述了概念（即自然）、自由和法之间的这种联系。第一次柏林法哲学讲座在1818-1819年、《法哲学原理》出版之前。其中没有收入正式出版文本中的第3节说道：“法的原则不在于自然，既不在外在自然，也不在人的主观自然，因此时人的意志是自然地被规定的，也就是说，是欲望、冲动和偏好的领域。法的领域是自由的领域，在其中，自由外在化自身并赋予自身以实存，此时尽管自然会出场，但却是作为一种没有独立性的东西[而出场的]。”\footnote{《自然法和国家法》(Natur- und Staatsrecht)，G. 霍迈尔根据黑格尔教授先生1818-1819年冬季学期演讲所作笔记，四开本德语手稿第155号，第9页。这一笔记以及下面引用的笔记藏于普鲁士文化遗产基金会柏林国家图书馆。[本段引文见《黑格尔全集》（历史考订版）第26卷第1分册，第239页。——译者]}正如法哲学所述，从抽象法的现象形式开始，直到家庭、市民社会和国家为止，自由的各种外化是人的历史—社会的（“客观的”）现实性的所有环节。毫无疑问，在所有这些环节中，处处都有“自然性”关系的身影。但问题在于，这些“自然性”关系是否从自身产生出一种支配那些形式的法律。在黑格尔看来，必然性——比如说，在与个体存在的关系中，国家的定在所具有的那种必然性——不再意味着，对于个体而言，必须在国家中生活乃是一种自然的规律。毋宁说，国家的必然性建立在自由自身所立的法律基础上。黑格尔对第3节的解说毫不含糊地指出：“这一节是由自然法之名引发的。自由与必然性的统一并不是通过自然，而是通过自由产生出来的。自然事物一成不变，自己并不能从规律中摆脱出来，以便自己制定法律。然而，精神则使自身从自然中挣脱出来，并且自己产生它的自然、它的法则本身。因此，自然不是法的生命。”与费希特的看法一样，黑格尔也认为自然法的名称“不过是习传下来的”，按照事物[本身]、法的“概念”，这一名称甚至是不正确的，因为“自然可以理解为：(1)本质、概念；(2)无意识的自然（真正的意义）”；“真正的名称”必须是“哲学法学”。\footnote{《自然法和国家法》，第9页。[本段引文见《黑格尔全集》（历史考订版）第26卷第1分册，第240页。——译者]黑格尔也在《世界历史哲学讲演录》的导论中指出了“自然”的双重含义。参见《历史中的理性》，第5版，J. 荷夫迈斯特编，汉堡，1955年，第117页：“如果‘自然’这个词描述的是一种事物的本质、概念，那么存在自然状态，自然法就是此种状态。这种法，按照其概念、按照精神的概念，应当归属于人。但这绝不能与精神在其自然性状态中之所是混淆在一起。”进一步参见《法哲学讲演录》，1824-1825年冬季学期，格里斯海姆笔记，四开本德语手稿545号，第5页：“[……]我们必须注意，‘自然’这一表述马上就具有一种歧义，一种重要的、会导致绝对错误的歧义。一方面，自然意味着自然性存在，正如我们发现自己在不同的方面是直接造成的那样，是我们存在的直接方面。与之相对并与之相区分，自然也是概念，事物的自然[本性]也就是事物的概念，就是事物的合乎理性方式的存在，这种事物完全不同于单纯自然性的东西。”[本段引文见《黑格尔全集》（历史考订版）第26卷第3分册，第1052页。——译者]}

我们已将黑格尔法思想发展的第三个阶段确定在1816年到1820年之间。这一阶段完全回到了《论自然法的科学探讨方式》一文想要通过“伦理自然”法的建构予以克服的那种自由哲学的基础上。在结束了对“哲学法学”基本概念的漫长理解过程（这一过程同时也澄清了哲学法学与传统自然法的历史关系）之后，现在对此也就无需赘言了。在黑格尔这里，自然与自由、自然规律与法权法则分离开来了。黑格尔在1822-1823年冬季学期的柏林法哲学讲座导言中说：“在自然中，有着一般说来一条规律存在的最高证明；而在法的法则中，事情之所以有效，并不是因为它存在，而是因为每个人都要求，它们应当符合一种自身的标准[……]如果我们对这两种法则[规律]的这种差别进行考察并追问，法的法则所归属的基地是什么，那么我们就会看到：法只产生于精神，因为自然没有法。”\footnote{《法哲学：按照黑格尔教授先生的讲演》(Philosophie des Rechts. Nach dem Vortrage des H. Prof. Hegel)，1822-1823冬季，霍托笔记，手写笔记第二册，第2页。在甘斯版的《法哲学》中，霍托所记黑格尔法哲学讲演导言部分的部分笔记作为序言的附释刊印（《全集》第8卷，第8页以下）。[本段引文见《黑格尔全集》（历史考订版）第26卷第2分册，第772页。——译者]}因而黑格尔决然地拒绝传统观念：就只有人可以认识和遵守自然规律而言，“自然规律”对人来说可以成为法的“范本”；人的历史性定在所唯一具有的法则，就是自由法则，不是“自然”，而是“概念”自身制定了这一法则。\footnote{参见上书，第3页：“对我们而言，事物的概念并不来自自然。”}但概念的自我立法——对于黑格尔而言，正是这种人的普遍法权能力的思想\footnote{参见上书，第5页：“这须被尊重为某种伟大之事：现在人因其为人，必须认为拥有权利，以至于他的人的存在要高于他的身份。”}，反抗现成的东西并打破了自然性的假象，在这种假象背后，“法的思想”、自由都还一直隐而不彰。卢梭、康德和费希特的观念论自然法首先把握到的这种思想将现成的法和对它的知识变为同一个东西。黑格尔在1824-1825年的讲座导言中说，尽管在以自然、自然需要和目的等为基础的“习惯的自然法”中，自由的思想并没有被有意排除掉，但“事实上，由于没有领会到这两种原则的独特性就接受了它们，从而怠慢了自由”。自由不能以自然的形式，而是必须以“概念”的形式来思维，在自由与其自身以及自由与自然之间，概念起着中介的作用。当自由以“概念”的要求出现时，这种要求就会主张：“自由仅仅指与法和伦理有关的自由，这样自然法之名就陷入了动摇之中。”\footnote{《法哲学讲演录》，格里斯海姆笔记，第9页以下。[本段引文见《黑格尔全集》（历史考订版）第26卷第3分册，第1054页。——译者]}

现在，黑格尔在这次讲座中明确地区分自然法和法哲学；讲座开篇这句话就包含了对《法哲学原理》这本书的“双重标题”的提示\footnote{参见上书，第3页。[见《黑格尔全集》（历史考订版）第26卷第3分册，第1051页。——译者]}：《自然法和国家学或法哲学原理》。我们的阐述表明，这个标题绝非偶然。它在一种非常具体的意义上保存了构成这部著作前史的那些环节。副标题“自然法和国家学纲要”涉及黑格尔法思想的出发点，以及他与古典政治学和现代自然法的争执。在这一过程中，黑格尔逐渐认识到两者的局限和克服两者的条件。这种理解的结果就是标题：法哲学，这一标题摆在整本书最前面，并且把自然法和国家学在自身中统一了起来。应当指出的是，由黑格尔引入的“哲学法”\footnote{黑格尔，《法哲学原理》，第3节。}的名称以及它所描述的事物，即自由的“概念”，在历史上产生于观念论自然法，在法哲学中，黑格尔知道他自己就是观念论自然法的继承人。